\solchapter{astro}

\sol{astro}{1} $d = 1/p = 10\un{pc} = 33$ light-years. For $10\%$ accuracy, the parallax must be ten times the error, $200\un{\mu as}$, giving $d = 5\un{kpc}$: a large part of the Milky Way. (That accuracy is reached only for bright stars; for faint ones the errors are several times larger.)

\sol{astro}{2} $0.1\times0.7\times(2\times10^{30})\times0.007\times c^2 = 9\times10^{43}\un{J}$, which at $3.8\times10^{26}\un{W}$ lasts $2.3\times10^{17}\un{s}$, or $7$ billion years. Detailed models give about $10$ billion.

\sol{astro}{3} $GM^2/RL\approx10^{15}\un{s}\approx30$ million years. Geologists and Darwin needed hundreds of millions of years for erosion and evolution; Kelvin insisted, on the authority of physics, that the Sun could not be that old. Physics was missing a source of energy: the nucleus, discovered decades later.

\sol{astro}{4} For a gas of $N = 2M/m_p$ particles, the virial theorem gives $3N\kB T = \abs U\approx\tfrac35GM^2/R$ (uniform sphere), so $\kB T\approx0.1\,GMm_p/R$: $T\approx2\times10^6\un{K}$. This is a mean; the centre is hotter, and detailed models give $1.5\times10^7\un{K}$. Even so, $\kB T\approx1\un{keV}$, a thousand times short of the $1\un{MeV}$ Coulomb barrier. Classically the Sun could not shine. Quantum tunnelling makes fusion possible for the fastest protons at a few times $\kB T$ (Problem~\ref{prob:gamow}).

\sol{astro}{5} $N = (R/\ell)^2 = (7\times10^8/10^{-3})^2\approx5\times10^{23}$ steps, a total path of $N\ell = 5\times10^{20}\un{m}$, which takes about $50\,000$ years. (The opacity actually varies hugely with depth, and detailed estimates range from about $10^4$ to $10^5$ years.) Neutrinos fly straight out in $2.3\un{s}$: the light we see today left the core in the Stone Age, but the neutrinos are fresh.

\sol{astro}{6} $t\propto M/L\propto M^{-2.5}$: $10\Msun$, $3\times10^7$ years; $3\Msun$, $6\times10^8$ years; $0.2\Msun$, $6\times10^{11}$ years, much longer than the age of the universe (every red dwarf ever formed is still shining). On Earth, single-celled life appeared within a billion years, but complex life took four. The $10\Msun$ star is hopeless; the $3\Msun$ star might at best host microbes. The red dwarf lives long enough, but its habitable zone is so close that planets there are tidally locked, and exposed to violent flares.

\sol{astro}{7} Setting $\sigma_TL/4\pi r^2c = GMm_p/r^2$, the $r$ cancels:
\[
  L_{\text{Edd}} = \frac{4\pi GMm_pc}{\sigma_T} = 3.3\times10^4\Lsun\,\frac{M}{\Msun}.
\]
For $10^9\Msun$, $L_{\text{Edd}} = 1.3\times10^{40}\un{W} = 3\times10^{13}\Lsun$, three hundred times a large galaxy. This is how a quasar outshines its entire host galaxy from a region the size of the solar system.

\sol{astro}{8} Equating $\hbar^2n^{5/3}/m_e$ with $GM^2/R^4$ and using $n\sim M/2m_pR^3$:
\[
  \frac{\hbar^2}{m_e}\left(\frac{M}{2m_p}\right)^{5/3}\frac{1}{R^5}\sim\frac{GM^2}{R^4},\qquad
  R\sim\frac{\hbar^2}{Gm_e(2m_p)^{5/3}}M^{-1/3}.
\]
For $M = \Msun$, $R\sim2000\un{km}$; a careful calculation gives about $6000\un{km}$, the size of the Earth. More mass compresses the star to a smaller radius, and at some point the electrons become relativistic (Problem~11).

\sol{astro}{9} For a circular orbit, $v^2/r = GM(r)/r^2$, so $M(r) = v^2r/G$: $9\times10^{10}\Msun$ at $8\un{kpc}$ and $5.6\times10^{11}\Msun$ at $50\un{kpc}$. The visible matter is only about $11\%$ of the mass inside $50\un{kpc}$: nearly $90\%$ is dark. Since $M\propto r$ and no edge is in sight, the dark halo extends farther still.

\sol{astro}{10} Jupiter: $(R_J/R_\odot)^2 = 1\%$. Earth: $(R_\oplus/R_\odot)^2 = 8.4\times10^{-5}$, about $84$ parts per million. The Earth's atmosphere makes stars twinkle at the $0.1\%$ level, so Earth-sized transits are hidden from the ground; they need space telescopes such as Kepler, which measured brightness to about $20$ parts per million. And an Earth-like orbit transits only once a year, so several years are needed to confirm one.

\sol{astro}{11} With $n\sim M/m_pR^3$, $P\sim\hbar c(M/m_p)^{4/3}R^{-4}$, while gravity requires $GM^2/R^4$. Both scale as $R^{-4}$, so the ratio of the two is independent of $R$: if gravity wins at one radius, it wins at every smaller radius too. They balance only for one particular mass, $\hbar c(M/m_p)^{4/3} = GM^2$, i.e.
\[
  M\sim\left(\frac{\hbar c}{G}\right)^{3/2}\frac{1}{m_p^2} = 3.7\times10^{30}\un{kg}\approx1.9\Msun.
\]
The careful calculation gives $1.4\Msun$. As a white dwarf gains mass, it shrinks, and its electrons are squeezed into higher momenta until they are relativistic. Then an electron's energy grows only as $p$, not as $p^2$, so squeezing raises the pressure only as fast as gravity's demand grows. Above the Chandrasekhar mass, gravity wins at every radius: the star collapses until electrons are forced into protons and it becomes a neutron star --- or, for an accreting white dwarf, carbon ignites and the star explodes as a type Ia supernova (Chapter~\ref{ch:cosmo}).

\sol{astro}{12} A contracting ball of gas heats up as the virial theorem requires, $\kB T\sim GMm_p/R$. But as $R$ shrinks, the electron degeneracy energy $\sim(\hbar^2/m_e)n^{2/3}\propto1/R^2$ grows faster, and eventually supports the star, which then cools. Balancing the two, the highest temperature reached is $\kB T_{\max}\sim G^2m_em_p^{8/3}M^{4/3}/\hbar^2$: it grows with mass. Setting it equal to the few million kelvin needed for hydrogen fusion gives a minimum mass of a few hundredths of a solar mass, within a factor of a few of the true $0.08\Msun$. At the top, the mass--luminosity relation $L\propto M^{3.5}$ grows faster than the Eddington luminosity $\propto M$; the Sun radiates at $1/33\,000$ of its Eddington limit, and the two would meet at $33\,000^{1/2.5}\approx60\Msun$. Near the limit, radiation pressure drives the outer layers off in winds, and the relation flattens, so the real cutoff is $100$--$150\Msun$. Between these two limits lies a factor of about $10^3$, set by fundamental constants: both can be traced to the Chandrasekhar combination $(\hbar c/G)^{3/2}/m_p^2\approx2\Msun$ of Problem~11, multiplied by modest powers of dimensionless numbers such as $\alpha$ and $m_e/m_p$.

\sol{astro}{13} Conservation of angular momentum, $I\omega\propto MR^2/P$, gives $P_{\text{ns}} = 3600\un{s}\times(12/6400)^2 = 13\un{ms}$. Conservation of flux, $BR^2$, gives $B = 100\un{T}\times(6400/12)^2 = 3\times10^{7}\un{T}$, close to typical pulsar fields. The equator flies off when $\omega^2R = GM/R^2$, so the shortest period is $2\pi\sqrt{R^3/GM} = 0.6\un{ms}$ for $1.4\Msun$, a spin of $1650$ revolutions per second. The fastest known pulsar spins at $716\un{Hz}$, less than half the break-up rate. (Real neutron stars bulge as they spin, lowering the limit somewhat.)
